% P10-Evaluationsfeinschliff, 16.09.2026: Aussagegrenzen und Anforderungsbilanz; keine neuen Messwerte.
% B-59, 15.09.2026: datierte Ledger-Nachauswertung, Stützung und Reichweite.
% Beleg: thesis-evidence/w2/ledger-repetitions-20260915/ (Originale unverändert).
% B-58, 15.09.2026: Umfang der bereits in D.2 belegten Stufen-Nachprüfung explizit.
% P10-4c-Nachtrag, 15.09.2026: Systemmodell der vier Hauptläufe explizit;
% an den archivierten config.json-Dateien und den F-Metriken geprüft.
% P10-4c, 14.09.2026: Ledger-Herleitung und Nachweisanschlüsse präzisiert.
% P10-4b, 14.09.2026: Begriffs-/Beleganschlüsse und abgegrenzte Prüfaussagen präzisiert.
% M19 (12.09.2026): Leserführung und präzise Ergebnisdarstellung;
% Messwerte und historische Nachweisgrenzen bleiben erhalten.
% M14 (11.09.2026): 7.2.2 verdichtet; Runden-Details in D.2/D.8,
% Verfahren, Pruefgrenzen und konsolidierter Stand bleiben im Haupttext.
% M02 (10.09.2026), B-16: Laufwahl gemäß Manifest; zeitliche
% Vorabfixierung nicht belegt. Einordnung und Beleggrenze in D.1.
% ============================================================
% §7.2 Untersuchungsdesign und Bewertungsverfahren — v02 (08.09.)
% NACHREVIEW: neue Bewertungsrunde 08.09. in 7.2.2 dokumentiert
% (2 Voll-Reviews → beauftragte KI-Adjudikation, 8 Änderungen,
% getrennt von der 8+3-Historie; Meeting-2-Konsistenz 9 IDs/0 Änd.);
% 7.2.3 + drei Zusatzanalysen (Typen/Evidenzinhalt/Szenario).
% Zusatz-HERKUNFT: runs/w2-nachreview-konsolidierung.json +
% kapitel-7/ADJUDIKATION-Auftrag-Kapitel-7-Nachreview.md + beide
% Review-Dateien (Review-Coverage-109-…/Review-F-Coverage-…).
% NEUE 8er-Struktur (Schreibplan v2.4 Zeile 7.2: Systemstand/
% Laufbestand · Fallauswahl ehrlich · Messgrenzen · Referenz+
% Urteilsregeln inkl. Zählregeln · Entstehungs-Disclosure).
% Ersetzt alt-7.2+7.3+7.4; Kollegen-Review 07.09. eingearbeitet:
% Einstiegssatz, 109/Verlauf, Gold-Precision-Korrektur, Nenner
% Referenzgültigkeit/Stützung, F1 raus, Prüfrunden getrennt,
% „vor der betreffenden Prüfung dokumentiert" statt vorregistriert.
% Budget ~2 S. + 1 Tabelle.
% ============================================================
% HERKUNFT: Verfahren/Grenzen = thesis-evidence/w2/
% messprotokoll-kapitel-7.md (inkl. Nachträge Gold v2 + Neube-
% rechnung + Autor-Volldurchsicht); Urteilsregeln =
% input/eval-labels/w2-matchregeln.md (Hash ① v1.2); Gold-
% Entstehung = input/eval-labels/w2-gold-aenderungslog-v2.md +
% benchmark-checksums.txt (Hash ①/②) + 2 Review-Dossiers 07.09.
% (kapitel-7/02-gold-pruefung/); Prüfrunden = runs/
% w2-official-manifest.txt (Chronik) + Z1-Verifikationsprotokoll-
% 2026-09-07.md; Läufe wie in §7.3/§7.5 einzeln benannt.
% ------------------------------------------------------------

\section{Untersuchungsdesign und Bewertungsverfahren}
\label{sec:eval-design}

Die quantitative Evaluation untersucht zwei bereits während der
Entwicklung verwendete Quellenfälle: ein kurzes Folge-Meeting mit
36 nummerierten Units und 31 Referenzaussagen sowie einen
zusammengeführten Interview- und Ausarbeitungsverlauf mit 136
Units und 109 Referenzaussagen. Die Fälle erfüllen
unterschiedliche Funktionen: Der Meeting-Fall ist kompakt und klar
strukturiert und dient als Treue-Fall; der Interview-Fall
durchläuft mehrere Gesprächs- und Ausarbeitungsphasen und dient
als Abdeckungs-Stressfall --- gut ein Drittel seines
Referenzbestands (37 von 109 Aussagen) sind späte
UI-Detailfestlegungen, was jede ungewichtete Abdeckungszahl
mitprägt. Dass beide Quellen entwicklungsbekannt waren, ist eine
ausgewiesene Grenze (Abschnitt~\ref{sec:eval-grenzen}); neue, unbekannte Quellen
standen für die Kampagne nicht zur Verfügung.

Die vier Hauptläufe --- je ein Lauf der Ledger- und der
F-Konfiguration pro Quellenfall --- verwendeten GPT-5.4 als
Systemmodell. Damit beruhen die Hauptvergleiche auf derselben
Modellwahl. Die ergänzende Modellreihe mit weiteren Modellen ist
gesondert im Anhang~\ref{anh:evaluation} ausgewiesen.

Die Evaluation verbindet retrospektive Auswertungen
dokumentierter Systemläufe mit ergänzenden, vor ihrer Ausführung
spezifizierten Kontrolltests (Abschnitt~\ref{sec:eval-kontrollen}); jeder Untersuchung
wird der tatsächlich verwendete Systemstand zugeordnet
(Zuordnung von Lauf-ID, Messpunkt und verfügbaren Stand- und
Modellangaben einschließlich dokumentierter Rekonstruktionslücken:
Evidenztabelle im Anhang~\ref{anh:evaluation}). Als
Benchmarklauf diente je Quellenfall und Konfiguration der
chronologisch erste vollständige Lauf gemäß der im Laufmanifest
dokumentierten Auswahlregel; deren Belegstand sowie Fehlversuche
der Kampagne sind im
Anhang~\ref{anh:evaluation} ausgewiesen. Beide Quellen lagen deterministisch vorsegmentiert
vor (nummerierte Units als gemeinsame Adressbasis aller
Konfigurationen); die Qualität dieser Segmentierung wird nicht
mitgemessen. Ein zusammenhängender finaler Integrationslauf wurde
nicht ausgeführt, entsprechend wird keine gemeinsame
Ausführbarkeit eines späteren Endstands behauptet. Die
Messgrenzen sind je Abschnitt gesetzt: Abschnitt~\ref{sec:eval-ledgerqualitaet} bewertet den
maschinellen Claim-Bestand \emph{vor} dem Human-Gate,
Abschnitt~\ref{sec:eval-kontrollen} die Kontrollen an der Speichernaht, Abschnitt~\ref{sec:eval-fortschreibung}
historisch belegte, teils menschlich mitgestaltete
Fortschreibungen.

\subsection{Referenzbestand}
\label{sec:eval-gold}

Die in Abschnitt~\ref{sec:evidenzgebundene-wissensgewinnung}
beschriebene offene Suche nach Auslassungen liefert für sich noch
keine feste Bezugsmenge für eine Abdeckungsquote. Deshalb legt ein
gesonderter Referenzbestand fest, welche Quellinhalte hier als
erwartet gelten; die von den Konfigurationen erzeugten Claims bilden
diesen Maßstab nicht selbst.

Der Referenzbestand entstand als KI-Erstentwurf, der vom Autor
Unit für Unit gegen die nummerierte Rohquelle konsolidiert wurde
--- ohne Einsicht in Ausgaben der bewerteten Konfigurationen.
Regeln (Annotationsleitfaden, Matchregeln) und Quellen wurden vor
der Validierung eingefroren (Hash-Stufe~1). Nach der
Hauptauswertung wurde der Bestand zwei separat durchgeführten,
quellenbasierten KI-Vollprüfungen unterzogen (jede Aussage vorwärts, jede
Unit rückwärts; gold-bewusst, nicht verblindet); die Befunde
wurden vom Autor einzeln adjudiziert, als versionierter Stand~v2
mit Änderungslog eingefroren (Hash-Stufe~2; im Interview-Fall kam
eine Aussage hinzu: 108\,$\rightarrow$\,109) und der Bestand vom
Autor vollständig durchgelesen (ohne Befund). Eine unabhängige
menschliche Zweitannotation fand nicht statt --- der Autor ist
zugleich Konsolidierer; dieser Prüfstand wird so berichtet, wie er
vorliegt. Alle nach der Revision betroffenen Urteile wurden gegen
v2 neu gefällt; Alt- und Neuwerte werden nicht vermischt.

\subsection{Urteilsverfahren und Prüfrunden}
\label{sec:eval-urteile}

Ein LLM-Evaluator beurteilt anhand der eingefrorenen Matchregeln,
ob eine Referenzaussage vollständig, teilweise oder nicht
enthalten ist. Mehrere Claims dürfen sie gemeinsam abdecken
(Bündelung); für die inhaltliche Korrektheit wird dagegen jeder
Claim einzeln beurteilt. Die systeminternen Prüfer ersetzen diese
externe Bewertung nicht.

Die Hauptauswertung vom 06.09. verband eine KI-gestützte
Urteilsbildung (Adjudikation) mit Beleg je Urteil, eine separate
KI-Nachprüfung (Verifikation) und eine menschliche
Stichprobenprüfung durch den Autor. Letztere umfasste
16 Urteilsblöcke des Stressfalls sowie die Prüfbasis der Lücken
ohne zugeordneten Prüfhinweis. Die nach der Gold-Revision am
07.09. neu gefällten Urteile erhielten keinen eigenen Verifikationsdurchgang.
Tabelle~\ref{t:anh-runden} dokumentiert die einzelnen Runden,
Prüfumfänge und Korrekturen; deren Zählungen bleiben getrennt.

Am 08.09. wurden beide Stressfall-Auswertungen vollständig
KI-nachgeprüft, einschließlich der bisherigen Volltreffer. Die
Prüfer kannten die bisherigen Urteile; eine Verblindung fand
nicht statt. Anschließend entschied eine gesonderte
KI-Adjudikation über die Änderungsvorschläge. Die dabei
präzisierten Regelauslegungen wurden auf beide Konfigurationen
angewandt und im Meeting-Fall gezielt an neun strukturgleichen
Fällen geprüft, ohne weitere Änderungen. Eine erneute
Vollprüfung des Meeting-Falls fand in dieser Runde nicht statt.
Die Stufenanalyse (Abschnitt~\ref{sec:eval-stufen}) erhielt
eine separate KI-Nachprüfung beider Statusänderungen und zehn
per Seed ausgewählter unveränderter Paare, insgesamt zwölf Paare;
auch sie war nicht verblindet und
wich von der vorgesehenen Lesereihenfolge ab. Eine unabhängige
menschliche Zweitannotation fand nicht statt.

Kapitel~\ref{chap:evaluation} berichtet den daraus entstandenen
\emph{konsolidierten Auswertungsstand}. Die Nachprüfung vom 08.09.
änderte Bewertungen derselben Systemausgaben, weder die Ausgaben
selbst noch den Referenzbestand. Frühere Werte bleiben im
Änderungsnachweis erhalten. Die grundlegenden Match- und
Zählregeln waren vor der Hauptauswertung eingefroren; spätere
Prüfrunden dokumentierten ihre Prüfregeln vor der jeweiligen
Prüfung. Die Nachreview-Runde enthält zusätzlich nachträgliche,
ergebnisbekannte Auslegungsentscheidungen. Ihre Begründungen und
die Auswirkungen alternativer Auslegungen auf die Kennzahlen
sind in Anhang~\ref{anh:evaluation},
\hyperref[anh:coverage-nachreview]{Abschnitt~D.8}, ausgewiesen.
Auch die Darstellung dieses Kapitels wurde nach der
Hauptauswertung überarbeitet.

Am 15.09. wurden zwei weitere archivierte Ledger-Ausgaben des
Stressfalls nach den bisherigen Regeln und dokumentierten
Auslegungen bewertet: je 109 Abdeckungsurteile sowie Korrektheit
und Quellenstützung aller Claims. Ein vor diesen Urteilen
festgehaltener Leitfaden bestimmte das Vorgehen. Ein dabei erkannter
unvollständiger Quellenbezug führte zur dokumentierten Erweiterung um eine
erneute Claim-Prüfung des Hauptlaufs; dessen Abdeckungsurteile blieben
unverändert. Insgesamt wurden 135 Claims KI-gestützt nachgeprüft.
Der Autor beurteilte drei ausgewählte Fälle zur Verbindlichkeit der Aussagen anhand von
Quellen-/Claim-Gegenüberstellungen. Diese fallbezogene Prüfung ist
keine unabhängige Zweitannotation. Verfahren, Revisionen und
Auslegungsalternativen stehen in
\hyperref[anh:ledger-wiederholungen]{Anhang~D.11}.

\subsection{Metriken und Zählregeln}
\label{sec:eval-metriken}

Tabelle~\ref{t:eval-metriken-v2} definiert die verwendeten
Kennzahlen nach dem tatsächlich angewandten Rechenverfahren.
$G$ bezeichnet den Referenzbestand, $|G|$ seine Anzahl an Aussagen.
Die Zähler $N_{\mathrm{full}}$, $N_{\mathrm{partial}}$ und
$N_{\mathrm{none}}$ erfassen vollständig, teilweise beziehungsweise
nicht abgedeckte Referenzaussagen. Als \emph{Lücke} zählt hier jede
nicht vollständig abgedeckte Aussage, also sowohl \emph{partial}
als auch \emph{none}.

% ============================================================
% Tabelle 7.2-1: Metrikdefinitionen (korrigiert) — v01 (07.09.)
% In Overleaf ablegen als: tables/eval-metriken-v2.tex
% Aufruf im Fließtext: \input{\tabledir/eval-metriken-v2.tex}
% Label: t:eval-metriken-v2
% Stil: chap6.1-Konvention. ERSETZT tables/eval-metriken.tex
% (alte Fassung: F1 [gestrichen laut Plan v2.1] + Precision-/
% Nenner-Definitionen wichen vom Rechenverfahren ab —
% Kollegen-Review 07.09., Punkte 3–5).
% HERKUNFT: Verfahren = w2-matchregeln.md §§2–7 + tatsächliche
% Berechnung in runs/w2-goldv2-eval-*.json (gold_escapes separat,
% supported nicht als Fehler); Zuordnungsregel = z1-stufenanalyse.md
% §4; Silent-Miss = Hauptmetrik (Schreibplan v2.1-Regel).
% ============================================================

\begin{table}[htbp]
  \centering
  \small
  \begin{tabular}{>{\raggedright\arraybackslash}p{0.26\textwidth}>{\raggedright\arraybackslash}p{0.64\textwidth}}
    \hline
    \textbf{Metrik} & \textbf{Definition (tatsächliches Rechenverfahren)} \\
    \hline
    \multicolumn{2}{l}{\emph{Ergebnisqualität und Quellenbindung}} \\
    Full Coverage & $N_{\mathrm{full}}/|G|$: Anteil der Referenzaussagen, deren Inhalt vollständig im Claim-Bestand enthalten ist (Bündelung mehrerer Claims erlaubt) \\
    Touched Coverage & $(N_{\mathrm{full}}+N_{\mathrm{partial}})/|G|$: Anteil vollständig oder teilweise enthaltener Referenzaussagen \\
    Inhaltliche Korrektheit & präzise Lesart: Anteil der Claims ohne \emph{festgestellte} Quellenverfälschung. Claims ohne quellenwidrigen oder verfälschenden Inhalt geteilt durch alle gelieferten Claims; geprüft gegen die gesamte nummerierte Quelle (nicht nur die zitierten Units). Die claimbezogene Quote hängt am Zuschnitt der Ausgabe (Bündelung verändert den Nenner) und ist zwischen verschieden granularen Beständen nur eingeschränkt vergleichbar. Gold-fremde, quellengetragene Inhalte zählen nicht als Fehler, sondern werden --- soweit erhoben (Anhang) --- gesondert als \emph{gold\_escapes} ausgewiesen (keine Gold-Deckungs-Präzision; Abweichung vom Matchregeln-P2 im Text dokumentiert). Abgrenzung: Semantische Stützung prüft die Belegqualität der \emph{zitierten} Units, diese Kennzahl die Verfälschungsfreiheit gegenüber der Quelle \\
    Referenzgültigkeit & gültige \emph{Einzelreferenzen} geteilt durch alle abgegebenen Einzelreferenzen (technische Auflösbarkeit der Unit-IDs) \\
    Semantische Stützung & direkt oder inferenziell gestützte Claims geteilt durch Claims mit mindestens einer auflösbaren Referenz; beurteilt wird die Vereinigung der referenzierten Units, unabhängig vom Selbst-Etikett des Systems \\
    \multicolumn{2}{l}{\emph{Rechenschaft (Umgang mit eigenen Lücken) --- Hauptmetrik: Silent-Miss}} \\
    Silent-Miss-Rate & $N_{\mathrm{ohne\ Hinweis}}/(N_{\mathrm{partial}}+N_{\mathrm{none}})$: Anteil unvollständig abgedeckter Referenzaussagen ohne zuordenbaren Prüfhinweis; das Gegenstück ($1-$ Silent-Miss) ist eine nachträgliche Hinweiszuordnungsquote, keine gemessene Wahrnehmung \\
    Adressierbarkeit & (voll gedeckte $+$ hinweis-zugeordnete unvollständige Referenzaussagen)$\,/\,|G|$ --- ergänzende Größe: Prüfpotenzial, keine erreichte Abdeckung, keine gemessene Behebung \\
    Signalbezogene Präzision & lückenrelevante Hinweise geteilt durch alle erzeugten Hinweise; ein zugeordneter Hinweis ist keine explizite Fehlermeldung (Prüfergebnisse werden mitberichtet) \\
    \hline
  \end{tabular}
  \caption[Definitionen der verwendeten Kennzahlen]{Definitionen der verwendeten Kennzahlen nach dem
    tatsächlich angewandten Rechenverfahren. Qualitäts- und
    Quellenbindungsmetriken bewerten den gelieferten Bestand;
    Rechenschaftsmetriken die nachträgliche Zuordenbarkeit
    eigener Auslassungen zu erzeugten Hinweisen.
    Prüfhinweise werden nie als erreichte Abdeckung
    gutgeschrieben.}
  \label{t:eval-metriken-v2}
\end{table}


\emph{Zuordnung und Zählung.} Ein Prüfhinweis darf mehrere Lücken
adressieren; jede Lücke und jeder Hinweis zählt je Kennzahl nur einmal.
Eine Lücke gilt als zugeordnet, wenn Hinweis und Referenzaussage
mindestens eine Quell-Unit gemeinsam haben. Einzeln dokumentierte
semantische Zuordnungen ergänzen diese Regel. Kennzahlen mit Nenner
null werden als nicht definiert ausgewiesen. Eine gültige Quellen-ID
belegt zunächst nur die technische Auflösbarkeit des Verweises.

% P10-6: Definition der tatsächlich verwendeten Korrektheitskennzahl zusammenhalten.
\begin{samepage}
\emph{Abweichung vom Messprotokoll.} Die eingefrorenen Matchregeln
sehen eine strikte \emph{Gold-Precision} vor: den Anteil der einzeln
vollständig gold-gematchten Claims an allen gelieferten Claims.
Quellengetragene Inhalte außerhalb des Referenzbestands gelten dort
nicht als Treffer. Diese Kennzahl wurde nicht berechnet.
Tatsächlich angewandt und hier berichtet wird die \emph{inhaltliche
Korrektheit}: der Anteil der Claims ohne festgestellten Widerspruch
zur Quelle oder Quellenverfälschung. Die Abweichung wurde für beide
Konfigurationen identisch angewandt und im Messprotokoll nachgetragen;
das eingefrorene Regeldokument bleibt unverändert erhalten.
\par
\end{samepage}


Quellengetragene Inhalte außerhalb des Referenzbestands zählen bei
diesem Verfahren nicht als Fehler. Eine nach der Gold-Revision
geprüfte Zählung solcher \emph{gold\_escapes} liegt nur für den
F-Treue-Fall vor (drei, alle gestützt). Das historische
F-Stressfall-Blatt enthält einen Zähler von 17 auf dem früheren
Referenzstand. Er wurde nach der Revision nicht neu geprüft und wird
nicht übernommen (Anhang~\ref{anh:evaluation},
\hyperref[anh:referenzpruefung]{Abschnitt~D.3}).

\emph{Gegenstand der Inhaltsbewertung.} Bewertet wird der
Aussagentext der Claims. Die fachliche Güte weiterer strukturierter
Felder, etwa Facetten oder Status, gehört nicht zu dieser Kennzahl;
systeminterne Prüferurteile ersetzen das externe Urteil nicht.
Plausible, aber in der Quelle nicht gestützte Ergänzungen gelten
nicht als quellenwidrig. Sie werden über die gesonderte Kennzahl der
\emph{semantischen Stützung} erfasst. Auch ein inhaltlich korrekter
Claim kann durch seine angegebenen Units nur teilweise gestützt sein,
wenn nötiger Gesprächskontext außerhalb dieser Auswahl liegt.
Bei F bewertete die bestehende Auswertung genau
eine redaktionelle Zuordnung als nicht vollständig quellengetragen:
die Zuweisung der zweiten Versionsnummer zu Dart, während die Quelle
zweimal Flutter nennt. Daraus ergibt sich die Stützung von 0{,}991.
Die ergänzende Claim-Prüfung vom 15.09. beschränkte sich auf den
Ledger; die Stützungsurteile von F wurden dabei nicht erneut geprüft.

\emph{Messgrenze.} Alle Inhaltsmetriken werden je Fall berechnet,
ohne Aggregation über die Fälle. Läufe ohne gültiges Artefakt gelten
als Ausfall und erhalten keine künstlichen Nullwerte. Der Aufwand
umfasst ausschließlich Input- und Output-Tokens desselben
Verarbeitungsausschnitts, einschließlich tatsächlich zugehöriger
interner Reparaturen und Wiederholungen. Zusätzliche fehlgeschlagene
Versuche der Kampagne werden gesondert ausgewiesen; die externe
LLM-Adjudikation zählt nicht zum Systemaufwand.

Die Trennung von Korrektheits- und Abdeckungsprüfung greift eine
Unterscheidung der aktuellen Evaluationsliteratur zu langen
generierten Texten auf: Samarinas et al.\ unterscheiden die
faktische Stützung generierter Aussagen von der Abdeckung
erwarteter Aspekte
% K6-TODO (Autor, Overleaf-Bib): [96] von "arXiv Version 4" auf die
% begutachtete Publikationsfassung umstellen — VERIFIZIERT 08.09. an
% aclanthology.org/2025.findings-acl.693: Samarinas, Krubner, Salemi,
% Kim, Zamani: "Beyond Factual Accuracy: …", Findings of the
% Association for Computational Linguistics: ACL 2025, S. 13468–13482,
% Wien, Juli 2025, DOI 10.18653/v1/2025.findings-acl.693.
% Fertiger BibTeX-Eintrag im NACHREVIEW-UMSETZUNGSPROTOKOLL.md.
\cite{samarinas2025coverage}. Die vorliegende Arbeit übernimmt
diese Trennung der Bewertungsrichtungen, operationalisiert
Abdeckung jedoch als ungewichtete Erhaltung festgelegter
Referenzaussagen --- jede zählt gleich; ein kombinierter
Gesamtscore im Sinne dieses Verfahrens wird nicht berechnet.

Die Nachreview-Runde ergänzt drei klar getrennte
\emph{Zusatzanalysen} ohne neue Gesamtmetrik: eine qualitative
Typisierung der Lücken beider Konfigurationen nach dem fehlenden
Bestandteil (Häufigkeiten, keine Wichtigkeitsskala); eine
Betrachtung, ob im Aussagentext eines Claims fehlender Inhalt in
\emph{tatsächlich gespeicherten} Belegzitaten erhalten ist (ein
anderer Artefaktumfang als die Claim-Abdeckung --- der lediglich
referenzierbare Ursprungstext zählt dort nicht); und eine
ausdrücklich hypothetische Szenariorechnung zu den
Queue-Zuordnungen (Abschnitt~\ref{sec:eval-luecken}). Diese Größen werden
nebeneinander berichtet und nicht addiert; Definitionen und
Herleitungen stehen im Anhang~\ref{anh:evaluation}.
